% --- INSTITUTIONAL TEMPLATE FED v7.4.9 ---
% Estándar de Reporteo APA 7 (Babel-Free)
\documentclass[12pt,spanish]{article}

\usepackage{graphicx}
\usepackage{geometry}
\usepackage{indentfirst} 
\usepackage{caption}
\usepackage{floatrow} 

% Forzar estilos de caption antes de cualquier carga
\floatsetup[figure]{capposition=top}
\floatsetup[table]{capposition=top}

\usepackage{booktabs}
\usepackage{float}
\usepackage{longtable}
\usepackage{threeparttable}
\usepackage{array}
\usepackage{calc}
\usepackage{amsmath, amssymb}
\usepackage{fancyhdr}
\usepackage{mathptmx} 
\usepackage[T1]{fontenc}
\usepackage{setspace} 
\usepackage{ragged2e} 
\usepackage{titlesec}
\usepackage{url}
\def\UrlBreaks{\do\/\do-\do\.\do\_\do\?\do\=\do\&\do\%}
\usepackage[hidelinks,breaklinks]{hyperref}

% --- Pandoc Compatibility Block (Iron Clad) ---
\providecommand{\tightlist}{%
  \setlength{\itemsep}{0pt}\setlength{\parskip}{0pt}}

\usepackage{xcolor}
\usepackage{fancyvrb}
\usepackage{framed}

% --- Code Highlighting (Pandoc Standard) ---
\AtBeginDocument{
  \definecolor{shadecolor}{RGB}{248,248,248}
}

% Compatibilidad Pandoc 3.x
\ifdefined\pandocbounded\else
  \newcommand{\pandocbounded}[1]{#1}
\fi

% --- Pandoc Citeproc (CSL) compatibility ---
\newlength{\cslhangindent}
\setlength{\cslhangindent}{1.27cm} 
\newlength{\csllabelwidth}
\setlength{\csllabelwidth}{3em}
\newcommand{\citeproctext}{} 
\newcommand{\citeprocitem}[2]{\item} 
\newenvironment{CSLReferences}[2] 
 {\begin{list}{}{
  \setlength{\leftmargin}{\cslhangindent}
  \setlength{\parsep}{0pt}
  \setlength{\itemsep}{6pt}
  \ifodd #1
   \setlength{\leftmargin}{\leftmargin + \csllabelwidth}
   \setlength{\itemindent}{-\csllabelwidth}
  \fi
  }
  \renewcommand{\bibitem}[2][]{\item} % Elimina los corchetes [ ]
 }
 {\end{list}}
\newcommand{\CSLBlock}[1]{#1\hfill\break}
\newcommand{\CSLLeftMargin}[1]{\parbox[t]{\csllabelwidth}{#1}}
\newcommand{\CSLRightInline}[1]{\parbox[t]{\linewidth - \csllabelwidth}{#1}\break}
\newcommand{\CSLIndent}[1]{\hspace{\cslhangindent}#1}

\makeatletter
\@ifundefined{paragraph}{\newcounter{paragraph}}{}
\@ifundefined{subparagraph}{\newcounter{subparagraph}}{}
\makeatother

% Solución al error 'No counter none defined' en longtable
\makeatletter
\newcounter{none}
\@ifpackageloaded{caption}{
  \DeclareCaptionType{none}
  \captionsetup[none]{name=Tabla}
}{}
\makeatother

\hypersetup{
  colorlinks=true,
  linkcolor=blue,
  filecolor=blue,
  citecolor=blue,
  urlcolor=blue
}

% --- Macros Institucionales ---
\newcommand{\fuente}[1]{
  \vspace{4pt}
  \begin{flushleft}
    \small
    \textit{Nota.} #1
  \end{flushleft}
}

% --- Localización Manual Determinista (Sin Babel) ---
\AtBeginDocument{
  \renewcommand{\tablename}{Tabla}
  \renewcommand{\figurename}{Figura}
  \renewcommand{\contentsname}{Índice}
  \renewcommand{\listtablename}{Índice de Tablas}
  \renewcommand{\listfigurename}{Índice de Figuras}
  \renewcommand{\abstractname}{Resumen}
  \renewcommand{\refname}{Referencias}
  
  \RaggedRight
  \setlength{\parindent}{1.27cm}
  \setlength{\parskip}{0pt}
}

\DeclareCaptionLabelSeparator{newline}{\\ \vspace{2pt}}
\captionsetup{
  position=top,
  justification=RaggedRight,
  singlelinecheck=false,
  labelfont=bf,
  textfont=it,
  labelsep=newline,
  font=small,
  skip=10pt
}
\captionsetup[figure]{name=Figura}
\captionsetup[table]{name=Tabla}

% Márgenes APA 7
\geometry{a4paper,margin=25.4mm,headheight=15mm}

% --- Encabezados ---
\pagestyle{fancy}
\fancyhf{} 
\fancyhead[L]{\includegraphics[height=1.2cm]{\logoUNL}}
\fancyhead[R]{\includegraphics[height=1.2cm]{\logoECONOMIA}}
\renewcommand{\headrulewidth}{0pt}
\fancyfoot[R]{\thepage}

\fancypagestyle{plain}{
  \fancyhf{}
  \fancyhead[L]{\includegraphics[height=1.2cm]{\logoUNL}}
  \fancyhead[R]{\includegraphics[height=1.2cm]{\logoECONOMIA}}
  \fancyfoot[R]{\thepage}
  \renewcommand{\headrulewidth}{0pt}
}

% --- Títulos ---
\titleformat{\section}{\normalfont\normalsize\bfseries\centering}{\thesection}{1em}{}
\titleformat{\subsection}{\normalfont\normalsize\bfseries}{\thesubsection}{1em}{}
\titleformat{\subsubsection}{\normalfont\normalsize\bfseries\itshape}{\thesubsubsection}{1em}{}
\titleformat{\paragraph}{\normalfont\normalsize\bfseries}{\theparagraph}{1em}{}
\titlespacing*{\section}{0pt}{12pt}{6pt}
\titlespacing*{\subsection}{0pt}{12pt}{6pt}
\titlespacing*{\subsubsection}{0pt}{12pt}{6pt}
\titlespacing*{\paragraph}{0pt}{12pt}{6pt}

\newenvironment{referenciasAPA}{
  \setlength{\parindent}{-1.27cm}
  \setlength{\leftskip}{1.27cm}
  \setlength{\parskip}{6pt}
}{\par}

\makeatletter\@ifpackageloaded{caption}{}{\usepackage{caption}}\AtBeginDocument{%
\ifdefined\contentsname
  \renewcommand*\contentsname{Table of contents}
\else
  \newcommand\contentsname{Table of contents}
\fi
\ifdefined\listfigurename
  \renewcommand*\listfigurename{List of Figures}
\else
  \newcommand\listfigurename{List of Figures}
\fi
\ifdefined\listtablename
  \renewcommand*\listtablename{List of Tables}
\else
  \newcommand\listtablename{List of Tables}
\fi
\ifdefined\figurename
  \renewcommand*\figurename{Figure}
\else
  \newcommand\figurename{Figure}
\fi
\ifdefined\tablename
  \renewcommand*\tablename{Table}
\else
  \newcommand\tablename{Table}
\fi
}\@ifpackageloaded{float}{}{\usepackage{float}}
\floatstyle{ruled}
\@ifundefined{c@chapter}{\newfloat{codelisting}{h}{lop}}{\newfloat{codelisting}{h}{lop}[chapter]}
\floatname{codelisting}{Listing}\newcommand*\listoflistings{\listof{codelisting}{List of Listings}}\makeatother\makeatletter\makeatother\makeatletter\@ifpackageloaded{caption}{}{\usepackage{caption}}
\@ifpackageloaded{subcaption}{}{\usepackage{subcaption}}\makeatother

\begin{document}

\begin{center}
    {\bfseries \Large UNIVERSIDAD NACIONAL DE LOJA}\\[0.3cm]
    {\bfseries \large FACULTAD JURÍDICA, SOCIAL Y
ADMINISTRATIVA}\\[0.2cm]
    {\bfseries CARRERA DE ECONOMÍA}\\[1.5cm]
    
    {\bfseries \large PROYECTOS DE INVERSIÓN PÚBLICA}\\[1cm]
    
    {\bfseries \Large Laboratorio de Proyectos de Inversión
Pública}\\[1.5cm]
\end{center}

\begin{center}
    \setstretch{1.1}
                        Erick Fabricio Condoy Seraquive \\
                \vspace{0.5cm}
    Econ. José Job Chamba \\ 
    19 de abril de 2026
\end{center}

\vspace{1cm}

\doublespacing
\setlength{\parindent}{1.27cm}

\section{Datos iniciales del
proyecto}\label{datos-iniciales-del-proyecto}

El ancho de texto en este reporte es \the\textwidth.

\subsection{Tipo de solicitud de
dictamen}\label{tipo-de-solicitud-de-dictamen}

Prioridad. La solicitud asegura el financiamiento vía Presupuesto
General del Estado para la mitigación efectiva de brechas en servicios
básicos, según la \textbf{Guía Metodológica para la Presentación de
Proyectos (Página 8)} (\textbf{guia\_planifica2021?}).

\subsection{Nombre Proyecto}\label{nombre-proyecto}

Proyecto de factibilidad para el mejoramiento del sistema de agua
potable en el barrio Carigán, parroquia Carigán del cantón Loja
provincia de Loja.

\subsection{Entidad UDAF}\label{entidad-udaf}

Gobierno Autónomo Descentralizado Municipal de Loja. Según el
\textbf{COOTAD}, esta entidad posee la competencia exclusiva para la
gestión de servicios de agua potable en su jurisdicción Asamblea
Nacional del Ecuador (2010).

\subsection{Entidad operativa
desconcentrada}\label{entidad-operativa-desconcentrada}

Unidad Municipal de Agua Potable y Alcantarillado (UMAPAL). Según el
\textbf{Reglamento Orgánico Funcional del GAD Loja}, UMAPAL es la unidad
técnica responsable de la operación y mantenimiento del sistema hídrico
cantonal.

\subsection{Gabinete Sectorial}\label{gabinete-sectorial}

Recursos Naturales, Hábitat e Infraestructura. Clasificado según la
estructura de coordinación definida en el \textbf{Decreto Ejecutivo
No.~91 (2024)} para proyectos de saneamiento ambiental.

\subsection{Sector, subsector y tipo de
inversión}\label{sector-subsector-y-tipo-de-inversiuxf3n}

De acuerdo con el catálogo de intervenciones definido en el
\textbf{Anexo 1 de la Guía de Planifica Ecuador}, el proyecto se
categoriza bajo los siguientes parámetros
(\textbf{guia\_planifica2021?}) detallados en la Tabla
\ref{tbl-sectores} y Tabla \ref{tbl-tipologias}:

\begin{longtable}{llll}

\caption{\label{tbl-sectores}Sectores y subsectores de intervención
definidos}

\tabularnewline

  \\
\toprule
\textbf{Macro sector} & \textbf{Sector} & \textbf{Código} & \textbf{Subsector} \\
\midrule
Social & Equipamiento urbano y vivienda & A0602 & Agua potable \\
\bottomrule

\end{longtable}

\fuente{Guía para la presentación de programas y proyectos de inversión pública.}

\begin{longtable}{llp{5cm}p{4cm}}

\caption{\label{tbl-tipologias}Tipologías de intervención definidas}

\tabularnewline

  \\
\toprule
\textbf{Cod.} & \textbf{Tipología} & \textbf{Conceptualización} & \textbf{Actividades relacionadas} \\
\midrule
T01 & Infraestructura & Procesos de adquisición, construcción, ampliación, mantenimiento, reparación, reposición, restauración de acervo físico. & Mejoramiento, reposición, construcción. \\
\bottomrule

\end{longtable}

\fuente{Guía para la presentación de programas y proyectos de inversión pública.}

\subsection{Plazo de ejecución}\label{plazo-de-ejecuciuxf3n}

18 meses. Este periodo se sustenta en la magnitud técnica de la obra
(12.5 km de red) y la complejidad topográfica del sector Carigán. El
cronograma contempla una fase de instalación y tendido de red de 14
meses, 2 meses destinados a pruebas de estanqueidad y desinfección, y 2
meses para el proceso de recepción provisional y liquidación técnica del
proyecto.

\subsection{Monto total}\label{monto-total}

\$ 1,050,000.00 USD. De acuerdo con el \textbf{Plan Plurianual de
Inversiones del PDOT Loja 2023-2027 (Página 1056)}, este monto se ajusta
a los techos presupuestarios definidos para la ampliación de redes en la
Zona 4 (\textbf{pdot\_loja\_2023?}).

\protect\phantomsection\label{refs}
\begin{CSLReferences}{1}{0}
\bibitem[\citeproctext]{ref-cootad2010}
Asamblea Nacional del Ecuador. (2010). \emph{Código orgánico de
organización territorial, autonomía y descentralización conocido como
COOTAD}. Registro Oficial.

\end{CSLReferences}

\end{document}
